\section{The composite hybrid for $\ell_1$-regularized PageRank}
\label{sec:composite-hybrid}

The previous sections use the 2026 RPPR paper mainly as evidence about
transient exploration.  We now address the regularized optimization problem
itself.  Catalyst applies directly to the composite objective
\cref{eq:rppr}; the nonsmooth term remains inside each proximal subproblem.
The hybrid is therefore well defined.  The first question--whether it at least
converges after a finite handoff--has a clean affirmative answer.

\subsection{Composite burn-in and proximal cleanup}

For
\[
  r_\rho(x):=\alpha\rho\norm{D^{1/2}x}_1,
  \qquad
  F_\rho(x)=f(x)+r_\rho(x),
\]
define the unit-step proximal map and fixed-point residual
\begin{equation}
  \mathcal T_\rho(x)
  :=\prox_{r_\rho}\bigl(x-\nabla f(x)\bigr),
  \qquad
  R_\rho(x):=x-\mathcal T_\rho(x),
  \label{eq:rppr-prox-map}
\end{equation}
where
\[
  \prox_{r_\rho}(v)
  :=\arg\min_z\left\{r_\rho(z)+\frac12\norm{z-v}_2^2\right\}.
\]
The map is weighted soft thresholding with coordinate threshold
$\alpha\rho\sqrt{d_i}$.

A composite Catalyst stage approximately minimizes
\begin{equation}
  H_t(z)
  :=F_\rho(z)+\frac\eta2\norm{z-y^{(t-1)}}_2^2,
  \qquad \eta=1-2\alpha.
  \label{eq:rppr-catalyst-subproblem}
\end{equation}
The shifted smooth Hessian is $Q+\eta I$ and has condition number two, exactly
as in the unregularized AESP subproblem.  Separability of $r_\rho$ makes
proximal-gradient and proximal-coordinate inner updates natural candidates.
This observation establishes algorithmic compatibility, but by itself does
not establish their degree-weighted local work from a signed Catalyst center.

\begin{algorithm}[t]
\caption{Proof-oriented composite hybrid for RPPR}
\label{alg:rppr-hybrid}
\begin{algorithmic}[1]
\REQUIRE $G,s,\alpha\in(0,1/2),\rho>0$ and a finite valid composite Catalyst
         burn-in.
\STATE Run any finite number of inexact Catalyst stages for
       \cref{eq:rppr-catalyst-subproblem}, obtaining $x^{(0)}$.
\STATE Optionally replace $x^{(0)}$ by its coordinatewise positive part.
\STATE Discard all outer momentum.
\FOR{$k=0,1,2,\ldots$}
  \IF{$\norm{D^{-1/2}R_\rho(x^{(k)})}_\infty
         \leq\alpha\eps_{\rm sol}$}
    \RETURN $x^{(k)}$.
  \ENDIF
  \STATE $x^{(k+1)}\leftarrow\mathcal T_\rho(x^{(k)})$.
\ENDFOR
\end{algorithmic}
\end{algorithm}

\subsection{Weighted contraction and a correct residual certificate}

\begin{lemma}[Weighted contraction of the RPPR proximal map]
\label{lem:rppr-weighted-contraction}
For every $x,y\in\R^n$,
\begin{align}
  \norm{D^{-1/2}(\mathcal T_\rho(x)-\mathcal T_\rho(y))}_\infty
  &\leq
  (1-\alpha)\norm{D^{-1/2}(x-y)}_\infty,
  \label{eq:rppr-infty-contraction}\\
  \norm{D^{1/2}(\mathcal T_\rho(x)-\mathcal T_\rho(y))}_1
  &\leq
  (1-\alpha)\norm{D^{1/2}(x-y)}_1.
  \label{eq:rppr-l1-contraction}
\end{align}
\end{lemma}

\begin{proof}
Weighted soft thresholding is coordinatewise nonexpansive.  Since
\[
  \mathcal T_\rho(x)
  =\prox_{r_\rho}\bigl((I-Q)x+b\bigr),
\]
we have componentwise
\[
  \abs{\mathcal T_\rho(x)-\mathcal T_\rho(y)}
  \leq (I-Q)\abs{x-y},
\]
because $I-Q=(1-\alpha)(I+W)/2$ is entrywise nonnegative.  Moreover,
\begin{align*}
  D^{-1/2}(I-Q)D^{1/2}
  &=\frac{1-\alpha}{2}(I+D^{-1}A),\\
  D^{1/2}(I-Q)D^{-1/2}
  &=\frac{1-\alpha}{2}(I+AD^{-1}).
\end{align*}
The first matrix has every row sum equal to $1-\alpha$, and the second has
every column sum equal to $1-\alpha$.  Their induced weighted
$\ell_\infty$ and $\ell_1$ norms are therefore $1-\alpha$.
\end{proof}

\begin{corollary}[Weighted proximal-residual certificate]
\label{cor:rppr-residual-certificate}
Let $x_\rho^\star$ be the unique minimizer of $F_\rho$.  Then
\begin{align}
  \norm{D^{-1/2}(x-x_\rho^\star)}_\infty
  &\leq
  \frac1\alpha\norm{D^{-1/2}R_\rho(x)}_\infty,
  \label{eq:rppr-infty-certificate}\\
  \norm{D^{1/2}(x-x_\rho^\star)}_1
  &\leq
  \frac1\alpha\norm{D^{1/2}R_\rho(x)}_1.
  \label{eq:rppr-l1-certificate}
\end{align}
In particular, the residual threshold in
Algorithm~\ref{alg:rppr-hybrid} certifies the requested weighted solution
error.  Omitting the factor $1/\alpha$ can stop up to that factor too early.
\end{corollary}

\begin{proof}
The minimizer is the unique fixed point
$x_\rho^\star=\mathcal T_\rho(x_\rho^\star)$.  In either weighted norm,
\[
\begin{split}
  \norm{x-x_\rho^\star}
  &\leq\norm{x-\mathcal T_\rho(x)}
       +\norm{\mathcal T_\rho(x)-\mathcal T_\rho(x_\rho^\star)}\\
  &\leq\norm{R_\rho(x)}+(1-\alpha)\norm{x-x_\rho^\star}.
\end{split}
\]
Rearrange.
\end{proof}

\begin{lemma}[Safe positive-part projection]
\label{lem:rppr-positive-part}
Let $x_+:=\max\{x,0\}$ coordinatewise.  Then
\begin{equation}
  F_\rho(x_+)\leq F_\rho(x).
  \label{eq:rppr-positive-part}
\end{equation}
\end{lemma}

\begin{proof}
Write $x=x_+-x_-$ with disjoint nonnegative parts.  The PageRank matrix $Q$
has nonpositive off-diagonal entries, so $x_+^\trans Qx_-\leq0$.  Therefore
\[
  x^\trans Qx
  =x_+^\trans Qx_++x_-^\trans Qx_- -2x_+^\trans Qx_-
  \geq x_+^\trans Qx_+.
\]
The linear coefficient $b$ is nonnegative, so removing $-x_-$ cannot increase
$-b^\trans x$, and the weighted $\ell_1$ term splits over $x_+$ and $x_-$.
\end{proof}

\begin{theorem}[Unconditional convergence of the composite hybrid]
\label{thm:rppr-hybrid-convergence}
Let any finite Phase-I method, including inexact composite Catalyst, return a
point.  Let $x^{(0)}$ be either that point or its nonnegative projection, and
run the proximal cleanup in Algorithm~\ref{alg:rppr-hybrid}.  Then
\begin{equation}
  \norm{D^{-1/2}(x^{(k)}-x_\rho^\star)}_\infty
  \leq
  (1-\alpha)^k
  \norm{D^{-1/2}(x^{(0)}-x_\rho^\star)}_\infty.
  \label{eq:rppr-hybrid-linear}
\end{equation}
Hence the composite hybrid converges after every finite Catalyst handoff.
\end{theorem}

\begin{proof}
Apply \cref{eq:rppr-infty-contraction} recursively.  The optional projection
is objective safe by \cref{lem:rppr-positive-part}; it is not required for the
contraction proof.
\end{proof}

\subsection{What this proves--and what remains open}

\Cref{thm:rppr-hybrid-convergence} settles correctness, not the COLT-style
local work target
\begin{equation}
  \softO\!\left(\frac1{\rho\sqrt\alpha}\right).
  \label{eq:rppr-target-work}
\end{equation}
A full proximal-gradient step may scan the support of the warm start and its
boundary.  The zero-initialized local ISTA proof uses coordinatewise
nonnegativity and support monotonicity, whereas a Catalyst point can overshoot
and need not lie below $x_\rho^\star$.  Thus iteration contraction alone does
not imply local degree work.

The most promising routes are now precise:
\begin{enumerate}[label=(\roman*),leftmargin=2.2em]
  \item certify a support envelope, reset to zero on that restricted region,
        and reuse the monotone local ISTA analysis;
  \item find a weighted proximal/KKT potential that decreases under local
        coordinate updates from an arbitrary warm start; or
  \item run a slightly over-regularized accelerated phase and use
        \cref{lem:kkt-charge} to charge only clearly inactive boundary nodes.
\end{enumerate}

\subsection{Assessment of the earlier internal draft}

The earlier draft supplied for this project already contained the essential
building blocks: the RPPR objective, weighted soft thresholding, local ISTA
and coordinate updates, the KKT conditions, the support-volume estimate
$\vol(S_\rho^\star)\leq1/\rho$, and a Catalyst construction.  Those ingredients
show that the present composite extension is not an artificial add-on.
The hybrid insight supplies the missing architecture: Catalyst need not remain
local through the final high-accuracy regime; it can be used only as a burn-in
before a momentum-free proximal cleanup.

Several points in the old draft should nevertheless be repaired before its
arguments are reused:
\begin{enumerate}[label=(\alph*),leftmargin=2em]
  \item the PageRank normalization and proximal step size must be fixed before
        comparing formulas from different sections;
  \item a proximal fixed-point residual is a solution-error certificate only
        after the $1/\alpha$ conversion in
        \cref{cor:rppr-residual-certificate};
  \item the statement that active sets converge to the optimal support needs a
        nondegeneracy/identification hypothesis and should not be used
        unconditionally;
  \item support monotonicity proved from the zero initialization does not
        transfer to a signed or overshooting Catalyst handoff; and
  \item standard composite Catalyst convergence does not by itself prove that
        its changing inner proximal problems have local degree-weighted work.
\end{enumerate}

Accordingly, the strongest defensible current RPPR statement is: the hybrid is
well defined and unconditionally convergent, while its graph-independent
accelerated local-work theorem remains a concrete open lemma rather than a
completed consequence of Catalyst.
